\chapter{Related Work}
\label{ch:related-work}

This chapter surveys the prior work on which the present thesis builds.
Section~\ref{sec:related-image} reviews image-based recognition systems
that classify static signs from individual video frames.
Section~\ref{sec:related-landmark} turns to landmark-based methods, which
replace pixel-level inputs with explicit hand and body keypoints.
Section~\ref{sec:related-temporal} surveys temporal models for
word-level recognition on the WLASL benchmark, covering recurrent,
three-dimensional convolutional, and transformer-based approaches.
Section~\ref{sec:related-mobile} discusses what is known about deploying
such models on mobile devices. Section~\ref{sec:related-position} closes
the chapter by situating the contribution of this thesis within the
surveyed landscape.

\section{Image-Based Sign Language Recognition}
\label{sec:related-image}

Early work on automatic sign language recognition often treated each
video frame as an independent image-classification problem, relying on
deep convolutional backbones pretrained on ImageNet to discriminate
static handshapes such as the letters of the American Sign Language
fingerspelling alphabet. The recent deep-learning survey of Rastgoo
et~al.~\cite{rastgoo2021slr} groups these approaches under the broader
umbrella of vision-based SLR and notes that, while convolutional
networks excel at frame-level appearance modelling, they cannot in
themselves capture the temporal dynamics that distinguish many
word-level signs.

The choice of backbone in this image-based line of work has tracked the
wider object-recognition literature. Convolutional families such as
ResNet remain common baselines, but recent work increasingly adopts
compound-scaled architectures. EfficientNet~\cite{tan2019efficientnet},
in particular, achieves competitive ImageNet accuracy with an
order-of-magnitude fewer parameters than earlier networks: its B0
variant reaches 76.3\% Top-1 with only 5.3 million parameters. This
favourable accuracy-versus-cost trade-off motivates its use as the
static-frame backbone of the present work
(Chapter~\ref{ch:implementation}).

Image-based pipelines remain useful for fingerspelling and isolated
static signs, but for word-level recognition---where motion, handshape
transitions, and two-handed coordination jointly carry meaning---they
must be combined with explicit temporal modelling. The next two
sections survey the two complementary strategies that have emerged:
replacing or augmenting the raw image with hand and body landmarks
(Section~\ref{sec:related-landmark}), and modelling temporal context
over the resulting frame-level representations
(Section~\ref{sec:related-temporal}).

\section{Landmark-Based Approaches}
\label{sec:related-landmark}

Image-based recognition is sensitive to background clutter,
illumination, signer appearance, and camera viewpoint. A complementary
line of work therefore reduces each frame to a small set of anatomical
keypoints---typically the joints of the hands and upper body---and
treats sign recognition as a problem over compact pose trajectories
rather than raw pixels. This shift has two practical consequences: the
input dimensionality drops by several orders of magnitude, and the
resulting representation factors out most appearance-related nuisance
variation, which in turn makes models more portable across signers and
recording conditions.

The de~facto extractor in this line of work is MediaPipe Hands of
Zhang et~al.~\cite{zhang2020mediapipe}, which delivers 21
three-dimensional landmarks per detected hand together with a
handedness label and a presence confidence score, in real time and on
commodity hardware. Its combination of accuracy, latency, and
on-device support has made it the default front end for many recent
sign-recognition pipelines, including the present work, which uses the
two-dimensional projections of these landmarks as its sole input
modality.

Pure landmark-only pipelines are not the only way to use pose
information. A complementary strategy is to keep the raw image but
use detected keypoints as a spatial guide for feature pooling. The
Fusion-3 method of Hosain et~al.~\cite{hosain2021fusion}, for example,
retains an Inflated 3D ConvNet RGB backbone and uses body-pose joints
to pool spatio-temporal feature maps from several depths of the
network, fusing the resulting pose-pooled classifiers with the
original 3D logits to reach 75.67\% Top-1 accuracy on WLASL-100. Such
hybrid systems achieve high accuracy at the cost of running a 3D
convolutional video backbone alongside the pose extractor; for the
mobile deployment target of this work that cost is prohibitive, which
motivates the landmark-only design choice revisited in
Section~\ref{sec:related-position}.

\section{Temporal Modeling for Word-Level Recognition}
\label{sec:related-temporal}

Word-level sign recognition is distinguished from fingerspelling by the
central role of motion: handshape transitions, two-handed coordination,
and trajectory shape jointly disambiguate signs that look identical in
any single frame. Models in this regime therefore have to integrate
information across an entire video clip, and the design space splits
cleanly into three families: recurrent networks built on top of
frame-level features, three-dimensional convolutional networks that
operate directly on video tensors, and more recent transformer-based
architectures that apply self-attention along the time axis.

The standard benchmark for evaluating these families is the Word-Level
American Sign Language dataset of Li et~al.~\cite{li2020wlasl}, which
provides four nested subsets of 100, 300, 1{,}000, and 2{,}000 sign
classes. The original WLASL paper accompanies the dataset with three
baselines that span the three families: a VGG-GRU pipeline that feeds
per-frame CNN features into a recurrent network, a Pose-TGCN that
applies a temporal graph-convolutional network over extracted body
pose joints, and an Inflated 3D ConvNet operating on the raw video.
On the smallest, hundred-class subset that anchors the comparisons in
this thesis (Chapter~\ref{ch:test-results}), these three baselines
reach 25.97\%, 55.43\%, and 65.89\% Top-1 accuracy respectively.

Recurrent networks remain a natural fit for variable-length sign
sequences. The Long Short-Term Memory unit of Hochreiter and
Schmidhuber~\cite{hochreiter1997lstm} introduced the constant-error
carousel and multiplicative input and output gates that allow gradient
information to flow over long time steps without vanishing, and
Schuster and Paliwal~\cite{schuster1997bidirectional} showed that
running such a recurrent unit in both temporal directions and
concatenating the resulting hidden states yields strictly more
contextual information per frame. The two-layer bidirectional LSTM
that this thesis runs on top of MediaPipe landmarks
(Chapter~\ref{ch:implementation}) is a direct instance of this
combination.

A second line of work pushes the recognition accuracy of
three-dimensional convolutional video backbones on WLASL by importing
additional supervision from outside the dataset itself. The TK-3D
ConvNet of Li et~al.~\cite{li2020transferring} retains the Inflated 3D
ConvNet backbone but augments it with weakly labelled sign clips
harvested from subtitled news broadcasts, then aligns the news domain
with the WLASL domain through a coarse classifier and a prototypical
memory of class centroids. The resulting model reaches 77.55\% Top-1
on WLASL-100, more than ten percentage points above the plain I3D
baseline. The improvement comes at the cost of a heavyweight 3D
convolutional pipeline plus a cross-domain training stage, both of
which assume RGB video input.

The most recent strand replaces both recurrence and three-dimensional
convolution with self-attention over pose sequences. The Sign
Pose-Based Transformer of Boh\'a\v{c}ek and
Hr\'uz~\cite{bohacek2022spoter} feeds body and hand keypoint
trajectories into a transformer encoder--decoder and reaches 63.18\%
Top-1 on WLASL-100 without ever touching a pixel of the original
video. This result demonstrates that, with a strong sequence model, a
pose-only pipeline can be competitive with mid-range RGB systems,
which is the regime in which the present work also sits.

\section{Mobile Deployment of Sign Language Models}
\label{sec:related-mobile}

Despite the steady accuracy gains reviewed above, very little of the
word-level sign-recognition literature reports on-device inference
benchmarks on the kind of consumer hardware that an accessibility tool
would actually run on. The reason is largely structural: the strongest
published systems on WLASL-100, namely TK-3D
ConvNet~\cite{li2020transferring} and
Fusion-3~\cite{hosain2021fusion}, both rely on Inflated 3D ConvNet
backbones whose runtime cost and memory footprint make them poorly
suited to commodity smartphones, even before adding the secondary pose
extractor that Fusion-3 requires or the cross-domain training pipeline
that TK-3D presumes. Landmark-only systems are, by construction, far
cheaper to evaluate at inference time---the input dimensionality is
bounded by a handful of keypoints rather than a full video
tensor---but the published landmark-based pipelines, including
SPOTER~\cite{bohacek2022spoter}, report accuracy on desktop hardware
rather than latency, energy, or model-size figures on a mobile
device. The recent survey of Rastgoo
et~al.~\cite{rastgoo2021slr} likewise treats deployment as a brief
practical note rather than as a primary evaluation axis.

This thesis takes the opposite emphasis. Because the target deployment
platform is a consumer Android smartphone
(Chapter~\ref{ch:implementation}), the recognition pipeline is
constrained from the outset to a lightweight on-device inference
runtime with a small model footprint, and the test-results chapter
(Chapter~\ref{ch:test-results}) reports model size and on-device
throughput alongside Top-1 accuracy on WLASL-100. We are not aware of
a directly comparable published on-device WLASL-100 benchmark on
commodity mobile hardware, which itself motivates the experimental
contribution of this work.

\section{Position of This Work}
\label{sec:related-position}

The present thesis occupies the corner of the design space that the
surveyed literature has explored least systematically: word-level
American Sign Language recognition driven entirely by MediaPipe hand
landmarks, evaluated on the WLASL-100 subset, and deployed end-to-end
on a consumer Android smartphone. Relative to the strongest reported
methods on this subset---the cross-domain TK-3D ConvNet of Li
et~al.~\cite{li2020transferring} and the pose-guided Fusion-3 of
Hosain et~al.~\cite{hosain2021fusion}---the design accepts a sizable
Top-1 accuracy gap (Chapter~\ref{ch:test-results}) in exchange for
discarding the RGB 3D convolutional pipeline that both systems
require. Relative to the pose-only transformer baseline of
Boh\'a\v{c}ek and Hr\'uz~\cite{bohacek2022spoter}, the pipeline
differs in two respects: the temporal model is a bidirectional LSTM
rather than a transformer encoder--decoder, and the evaluation
explicitly reports on-device latency and model size alongside Top-1
accuracy. The resulting contribution is therefore not a new accuracy
record on WLASL-100, but a demonstration that a landmark-only,
mobile-deployable pipeline can be evaluated on the same benchmark
under the same protocol, with all of its accuracy, footprint, and
latency trade-offs reported openly.
